%%
%% =====================================================================
%%  elsarticle 结构骨架
%% =====================================================================
%%  ⚠️  这不是 Elsevier 官方模板原件，而是按 elsarticle.cls 的**必需结构**
%%      写成的骨架，用途有两个：
%%        1. 作者手头没有下载官方模板时，也能生成结构正确、可直接编译的稿件；
%%        2. 作为「模板模式」的最小可用示例。
%%
%%      投正式期刊前，请用你从 Elsevier 官网下载的模板主文件替换它：
%%        python scripts/sci.py run 你的稿件 --journal elsevier-generic \
%%               --template /path/to/elsarticle-template-num.tex
%% =====================================================================
%%
\documentclass[preprint,12pt]{elsarticle}

\usepackage{amsmath}
\usepackage{amssymb}
\usepackage{graphicx}
\usepackage{booktabs}

\journal{[待作者填写：目标期刊名]}

\begin{document}

\begin{frontmatter}

\title{Template title placeholder}

\author[aff1]{Template Author}
\ead{template.author@example.com}
\affiliation[aff1]{organization={Template Institution}, city={City}, country={Country}}

\begin{abstract}
Template abstract placeholder.
\end{abstract}

\begin{keyword}
keyword one \sep keyword two \sep keyword three
\end{keyword}

\end{frontmatter}

\section{Introduction}
Template body placeholder.

\bibliographystyle{elsarticle-num}
\bibliography{refs}

\end{document}

\endinput
